%!TEX root = main.tex

We implement a solver  $\ostrichseq$ on top of the string solver OSTRICH~\cite{CHL+19} and the SMT solver $\princess$~\cite{princess08}, using about 5,000 lines of Scala code. The core functionality, preimage computation for dedicated string operations such as $\mywrite$, $\subseq$, $\elem$ and $\myseqlen$ is implemented in about 2,000 lines of code and is integrated directly into the OSTRICH framework. To support the representation of sequences, we built a lightweight library of automata which treat the separator as a special transition. The library is implemented in about 1,000 lines of code. 
%The remaining code is used for the integration with OSTRICH and for the implementation of the decision procedures for the theory of sequences.

%Moreover, we do experiments to evaluate the effectiveness of our approach.

%The experiment is conducted on a benchmark suite comprising the constraints generated from real-world and hand-crafted templates. We compare our implementation with the state-of-the-art SMT solvers, including Z3, and cvc5. The results show that our approach is efficient and capable of solving constraints that existing solvers cannot.


\subsection{Benchmarks}

To evaluate the effectiveness of $\ostrichseq$, we curate two benchmark suites, $\seqbase$ and $\seqext$,
where $\seqbase$ only uses operations that are directly supported by some existing SMT solvers while $\seqext$ uses some string sequence operations that are not directly supported by any existing SMT solvers.

\smallskip
\noindent\emph{$\seqbase$.}
The $\seqbase$ benchmark suite comprises $\arraylogic$ formulas that contain only the generic sequence operations, that is, sequence operations that can be applied to the elements of any type (not necessarily string), such as sequence length $\myseqlen$, sequence concatenation $\cdot$, subsequence $seqt[it_1, it_2]$, sequence read $\myat(seqt, it)$ and sequence write $seqt[it \rightarrow strt]$.
%
$\seqbase$ is designed to facilitate a fair comparison with existing solvers since they do not support the string-specific sequence operations. %where the sequences and strings are tightly related %(e.g. the constraints that contain $\mysplit_e$ and $\myjoin_u$ that convert between strings and sequences).
$\seqbase$ contains 140 $\arraylogic$ constraints that are generated from three templates. %Specifically,
\begin{itemize}
  \item  \polished{The first template $\svaru_1\in e_1 \wedge \svars_2 = \svars_1[n_1\rightarrow\svaru_1, n_2\rightarrow\svaru_1] \wedge \svaru_2 = \myat(\svars_2, m_1)\cdot\myat(\svars_2, m_2) \wedge \svaru_2 \in e_2 \wedge \myseqlen(\svars_2) < \mystrlen(\svaru_2)$ is utilized to curate 40 random  instances, where
  $\svaru_1$ is a string variable, $\svars_1,\svars_2$ are sequence variables, 
  %$s$ is a random generated sequence constant,
  and $e_1,e_2$ are randomly generated regular expressions.
  In the first 20 instances, $n_1, n_2, m_1, m_2$ are integer variables while in 
  the other 20 instances, they are random integers ranging from 0 to 10.}
  % \hide{Generally speaking, the template writes the same string to different positions of a sequence and then checks the length property and regular properties of elements.}
        %    
  \item \polished{The second template $\svaru_1\in e_1 \wedge \svars_2 = \svars_1[n_1\rightarrow\svaru_1, n_2\rightarrow\svaru_2] \wedge \svars_3 = \svars_2[n_3, len_1]\cdot\svars_2[n_4,len_2] \wedge \svaru_3 = \myat(\svars_3,m_1)\cdot\myat(\svars_3,m_2)\cdot\myat(\svars_3,m_3)\cdot\myat(\svars_3,m_4)\wedge\svaru_3\in e_2 \wedge \myseqlen(\svars_3) < \mystrlen(\svaru_3) + 1$ is utilized to curate 40 random instances, where $\svaru_1,\svaru_2$ are string variables, $\svars_1$, $\svars_2,\svars_3$ are sequence variables, and $e_1,e_2$ are randomly generated regular expressions.
  Similar to the first template,
  $n_i, m_i, len_1,len_2$ for $i\in [1,4]$  are integer variables in the first 20 instances but are random integers ranging from 0 to 10 in the other 20 instances.}
  %As above, we then replace all integer constants in this template with integer variables to produce an additional 20 instances.}
  % \hide{Briefly, the template generates a new sequence from an initial sequence by concatenating two overlapping subsequences and then checks the length and regular properties.}
  \item \polished{The third template $\svars_1 = \svars_0[m_1\rightarrow \svaru_1]\wedge\cdots\wedge\svars_{i} = \svars_{i-1}[m_i\rightarrow \svaru_i]\wedge\svaru_1=\myat(\svars_1, n_1)\wedge\cdots\wedge\svaru_j=\myat(\svars_j, n_j)\wedge\svaru_1\in e_1 \wedge\cdots\wedge \svaru_j\in e_j$ is utilized to curate 60 random instances, where $i$ and $j$ are randomly selected from the range $0$ to $20$, $\svars_0 \cdots \svars_i$ are sequence variables,  $\svaru_1\cdots \svaru_j$ are string variables, $e_1\cdots e_j$ are randomly generated regular expressions, and 
  $m_1,\cdots m_i$, $n_1\cdots n_j$ are randomly generated integers constant ranging from 0 to 5. In this template, write ($\svars[\cdot\rightarrow\cdot]$) and read ($\myat(\svars,\cdot)$) operations are performed randomly on a string sequence, and the strings read from the sequence are checked against some regular expressions.}
  % \item The third template $\svars_1 = \svars[m_1\rightarrow \svaru_1]\wedge\svars_2 = \svars_1[m_2\rightarrow \svaru_2]\wedge\cdots\wedge\svaru_1=\myat(\svars_1, n_1)\wedge\svaru_2=\myat(\svars_2, n_2)\wedge\cdots\wedge\svaru_1\in e_1 \wedge \svaru_2\in e_2 \wedge \cdots$ is utilized to generate 60 random instances, where $\svaru_1$, $\svaru_2$ are string variables, $\svars$, $\svars_1$, $\svars_2$ are sequence variables,  $e_1$, $e_2$ are randomly generated regular expressions, and 
  % $m_1$, $m_2$, $n_1$, and $n_2$ are randomly generated  integers constant. In this template, write ($s[\cdot\rightarrow\cdot]$) and read($\myat$) operations are performed randomly on a string sequence, and the strings read from the sequence are checked against some regular expressions.
\end{itemize}



%The benchmark contains 100 constraints generated from 3 different templates. The templates are aimed to cover the basic operations of string sequences. The templates and corresponding number of instances are as follows:

%Note that the regular properties of the sequence are defined by concatenating the first and the second string element in the sequence and checking whether the concatenated string satisfies a regular expression. This benchmark is designed to cover the basic operations of string sequences. 

\paragraph*{$\seqext$.}
The $\seqext$ benchmark suite comprises 60 $\arraylogic$ formulas that use specific string sequence operations, namely, $\myfilter_e$, $\mymatchall_e$, $\mysplit_e$ and $\myjoin_u$. Among them, 19 formulas are manually generated for the unit tests of $\ostrichseq$, and 41 formulas are generated from the real-world JavaScript programs collected from GitHub. Moreover, to compare with SMT solvers that do not support the considered string sequence operations, we rewrite them using the basic string operations whenever possible. ($\mymatchall_e$ is not rewritten since it cannot be expressed using existing string operations.) 
%We partition these formulas into two categories: $\withoutmatch$ and $\withmatch$. The former contains constraints that do not use $\mymatchall_e$, while %the $\withmatch$ category 
%the latter uses it. %contains the constraints that use the operation $\mymatchall_e$.  

%60 constraints generated by hand. 
%41 of them are transformed from the real-world JavaScript programs snippets, and 19 of them are collected from the unit tests of our solver. 
%The benchmark is designed to cover the extended operations of string sequences. 

%%%%%%%%%%%%%%%%%%%%%%%%%%%%%%%%%%%%%
\hide{
  Since we support some extended functionalities of string sequences such as \verb|str.split| and \verb|seq.join|, we divide the benchmark suite into two parts: \textbf{Seq-Base} and \textbf{Seq-Ext}. The \textbf{Seq-Base} benchmark contains the constraints with only the basic operations of string sequences, such as \textbf{length, concatenation, subsequence, read and write}. Note that we also use write in the benchmark, which can be transformed to concatenation, subsequence and length.
  The \textbf{Seq-Ext} benchmark contains the constraints with the extended operations of string sequences, such as \textbf{split, join, match and filter}.

  \begin{itemize}
    \item \textbf{Seq-Base}: The benchmark contains 100 constraints generated from 3 different templates. The templates are aimed to cover the basic operations of string sequences. The templates and corresponding number of instances are as follows:
          \begin{itemize}
            \item \texttt{Template 1}: 20 instances, which write the same string to different positions of the sequence and then check the regular properties of the sequence.

            \item \texttt{Template 2}: 20 instances, which generate a new sequence from the initial sequence by concatenating two overlapping subsequences and then checking the regular properties of the sequence.
            \item \texttt{Template 3}: 60 instances, which randomly read and write to a sequence of strings and then check the regular properties of the read strings.
          \end{itemize}
          Note that the regular properties of the sequence are defined by concatenating the first and the second string elements in the sequence and checking whether the concatenated string satisfies a regular expression. This benchmark is designed to cover the basic operations of string sequences.
    \item \textbf{Seq-Ext}: The benchmark contains 60 constraints generated by hand. 41 of them are transformed from real-world JavaScript program snippets, and 19 of them are collected from the unit tests of our solver. The benchmark is designed to cover the extended operations of string sequences.
  \end{itemize}
}
%%%%%%%%%%%%%%%%%%%%%%%%%%%%%%%%%%%%%

\subsection{Experimental results}
%We compare with cvc5, Z3, Z3-noodler, $\ostrich$ and $\princessarr$ on them.

We conduct experiments by comparing $\ostrichseq$ with SOTA solvers, including cvc5, Z3, Z3-noodler, $\ostrich$ and $\princessarr$ (referring to an array-theory-based solver implemented within $\princess$). $\seco$~\cite{JezLMR23} is excluded %in the comparison 
due to soundness issues. Additionally, \verb|foldleft| and \verb|map| operations are not used to simulate $\myjoin$, as most constraints return unknown under such usage in Z3.
%
On $\seqbase$, we evaluate $\ostrichseq$ along with sequence solvers cvc5, Z3 and $\princessarr$. For $\seqext$, we rewrite the constraints using only the basic string operations and compare $\ostrichseq$ with string solvers cvc5, Z3, Z3-noodler and $\ostrich$.

All experiments are run on a 24 $\times$ 3.1 GHz-core server with 185 GB RAM, where each solver is started as a single thread, with a 60 seconds time limit and without the memory limit.

\begin{table}[t]
\setlength{\tabcolsep}{6pt}
%\renewcommand{\arraystretch}{1.1}
%\vspace{-6mm}
  \centering
    \caption{\polished{Experimental results on $\seqbase$.}
  	%$\ostrichseq$ can solve all the constraints in $\seqbase$ and is more efficient than cvc5 and Z3. $\princessarr$ is slightly more efficient than $\ostrichseq$.
  } 
    \begin{tabular}{| c | c | c | c | c |}
      \hline
                      & cvc5        & Z3    & $\princessarr$ & $\ostrichseq$ \\
      \hline
      sat             & 48          & 14    & \textbf{52}    & \textbf{52}   \\
      \hline
      unsat           & 87          & 59    & 77             & \textbf{88}   \\
      \hline
      solved          & 135         & 73    & 129            & \textbf{140}  \\
      \hline
      unknown/timeout & 5           & 67    & 11             & \textbf{0}    \\
      \hline
      avg. time (s)   & 3.53        & 24.36 & 6.84           & \textbf{3.23}          \\
      \hline
    \end{tabular} %
%  \end{center}
  \label{tab:seq-base}
%\vspace{-8mm}
\end{table}


The results on $\seqbase$ are reported in Table~\ref{tab:seq-base}.
We can see that $\ostrichseq$ outperforms all the other solvers on $\seqbase$ in terms of both the number of solved instances and the solving time per instance. Specifically,   $\ostrichseq$ solves all the \polished{140} instances in $\seqbase$, while cvc5, Z3 and $\princessarr$ only solve \polished{135}, \polished{73} and \polished{129} instances, respectively. Moreover,  $\ostrichseq$ is more efficient than cvc5, Z3 and $\princessarr$ ($\ostrichseq$ spends \polished{3.23} seconds in solving each instance from $\seqbase$ on average, while cvc5, Z3 and $\princessarr$ spend \polished{3.25}, \polished{24.36}, and \polished{6.84} seconds, respectively). 

% \hide{However, $\ostrichseq$ is slightly slower than $\princessarr$, which solves all the instances in $\seqbase$ in 4.04 seconds on average. This is because $\princessarr$ utilizes the optimization techniques that are specialized to the array theory, making it more efficient for solving constraints that involve only read and write operations.}


The results on $\seqext$ are reported in Table~\ref{tab:seq-ext}. Since the SOTA sequence solvers do not support $\myfilter_e$, $\mymatchall_e$ or $\mysplit_e$ directly, the $\arraylogic$ constraints are rewritten to string constraints (instead of solving the $\arraylogic$ constraints directly) when we compare with them on $\seqext$.
% we fulfill the comparison with string solvers on $\seqext$ by 
%call string solvers use them to solve the string constraints generated from the $\arraylogic$ constraints (instead of solving the $\arraylogic$ constraints directly). 
Moreover, while we can use {\sf str.replace\_re\_all} to encode $\myfilter_e$ and $\mysplit_e$, the encoding of $\mymatchall_e$ requires finite-state transducers which have not been supported by most SOTA string solvers. %(except for {\sf OSTRICH}, \polished{but we do not implement this transformation in $\ostrich$, as it has already been applied in $\ostrichseq$}). 
%\denghang{{\sf OSTRICH} offer built-in support for finite-state transducers, but defining a transducer within an smtlib2.7 file is not feasible}). 
As a result, we do \emph{not} generate the string constraints for the 12 instances in $\seqext$ that contain $\mymatchall_e$, which are denoted by 
%We use 
$\seqext\mbox{\sf -M}$, whereas the rest 48 instances are %to denote the 12 instances in $\seqext$ that contain $\mymatchall_e$ and 
denoted by $\seqext\mbox{\sf -N}$. %to denote the other 48 instances in $\seqext$.


\begin{table}[t]
\setlength{\tabcolsep}{2pt}
%\renewcommand{\arraystretch}{1}
%\vspace{-6mm}
\caption{Experimental results on $\seqext$.
	% $\withmatch$ refers to the constraints that contain the operation $\mymatchall_e$, while $\withoutmatch$ refers to the constraints that do not contain the operation $\mymatchall_e$. The results show that $\ostrichseq$ can solve all the constraints in $\seqext$ and is more efficient than cvc5 and $\ostrich$. Z3 and Z3-noodler solve only a few constraints but are much faster than others.
}
%\vspace{-3mm}
\centering
    \begin{tabular}{| c| c| c | c | c | c| c|}
      \hline
                                              &                 & cvc5  & Z3   & Z3-noodler    & $\ostrich$ & $\ostrichseq$ \\
      \hline
      \multirow{3}{*}{$\seqext\mbox{\sf -N}$} & sat             & 19    & 4    & 4             & 20         & \textbf{23}   \\
      \cline{2-7}
                                              & unsat           & 21    & 6    & 4             & 18         & \textbf{25}   \\
      \cline{2-7}
                                              & solved          & 40    & 10   & 8             & 38         & \textbf{48}   \\
      \cline{2-7}
                                              & unknown/timeout & 8     & 38   & 40            & 10         & \textbf{0}    \\
      \cline{2-7}
                                              & avg. time (s)   & 10.83 & 0.13 & \textbf{0.05} & 15.92      & 3.53          \\
      \hline
      \multirow{3}{*}{$\seqext\mbox{\sf -M}$} & sat             & -     & -    & -             & -          & \textbf{12}   \\
      \cline{2-7}
                                              & unsat           & -     & -    & -             & -          & \textbf{0}    \\
      \cline{2-7}
                                              & solved          & -     & -    & -             & -          & \textbf{0}    \\
      \cline{2-7}
                                              & unknown/timeout & -     & -    & -             & -          & \textbf{0}    \\
      \cline{2-7}
                                              & avg. time (s)   & -     & -    & -             & -          & \textbf{1.77} \\
      \hline
    \end{tabular}
  \label{tab:seq-ext}
%\vspace{-8mm}
\end{table}


From Table~\ref{tab:seq-ext}, we can see that Z3 and Z3-noodler solve only 10 and 8 out of 48 constraints in $\seqext\mbox{\sf -N}$, respectively, and cvc5 and $\ostrich$ solve 40 and 38 instances, respectively,
indicating that the complexity of the string constraints generated from the $\arraylogic$ constraints. %are too complex to be solved by string solvers successfully in general. 
On the other hand, $\ostrichseq$ solves all of them, demonstrating its superiority in solving the $\arraylogic$ constraints with string sequence operations. For efficiency, %the solving time per instance, 
$\ostrichseq$ spends 3.53 seconds per instance on average, significantly less than cvc5 and $\ostrich$ (10.83 and 15.92 seconds on average, respectively). For the constraints in $\seqext\mbox{\sf -M}$, $\ostrichseq$ successfully solves all 12 instances with an average time of 1.77 seconds.
As aforementioned, for the constraints in $\seqext\mbox{\sf -M}$, we do not generate string constraints or compare against cvc5, Z3, Z3-noodler and $\ostrich$, since these constraints cannot be expressed using only the basic operations of string theory.
%are unable to solve any of these instances, as $\mymatchall_e$ cannot be rewritten using only the basic operations of string theory.
%In contrast, cvc5, Z3, Z3-noodler are unable to solve any of these instances, as $\mymatchall_e$ cannot be rewritten using only the basic operations of string theory.



%The results of the Seq-Base benchmark are shown in Table~\ref{tab:seq-base}. The results show that our solver is slower than Z3 and cvc5 in some cases, but it is still able to solve all the constraints. 

%%%%%%%%%%%%%%%%%%%%%%%%%%%%%%%%%%%%
%%%%%%%%%%%%%%%%%%%%%%%%%%%%%%%%%%%%
% \hide{
%   The results of the Seq-Ext benchmark are shown in Table~\ref{tab:seq-ext}. Since the extended operations are not supported by Z3 and cvc5, we rewrite the constraints to the basic operations of string theory, where we use string $'\#aa\#bb'$ to stand for the sequence $['aa', 'bb']$. \textbf{filter} and \textbf{match} are not rewrote since they are transducer which is not supported even in string theory. The results show that our solver OSTRICH outperforms Z3 and cvc5 in this benchmark, solving all the constraints with a competetive time. 28 unknown constraints of cvc5 and Z3 are due to filter and match. 17 unknown constraints of cvc5 and 25 unknown constraints of Z3 are due to the axioms we introduced to simulate the extended operations such as \textbf{split} and \textbf{join}.
% }
%%%%%%%%%%%%%%%%%%%%%%%%%%%%%%%%%%%%
%%%%%%%%%%%%%%%%%%%%%%%%%%%%%%%%%%%%


